\section{Tiền huấn luyện và lời nhắc}
\label{sec:ch02-tien-huan-luyen-va-loi-nhac}

Kiến trúc chỉ mô tả cách các token có thể tương tác; nó chưa cho biết mô hình
đã học những quy luật nào. Trong giai đoạn \tn{tiền huấn luyện}{pretraining},
một mô hình sinh tự hồi quy được tối ưu để dự đoán token tiếp theo trên một
tập văn bản lớn. Kết quả là cùng một kiến trúc có thể tạo ra hành vi rất khác
nhau khi trọng số được học từ dữ liệu khác nhau.

Sau tiền huấn luyện, \tn{lời nhắc}{prompt} trở thành một phần của điều kiện
đặt lên mô hình tại thời điểm dùng. GPT-3 khảo sát thiết lập few-shot, trong
đó mô tả nhiệm vụ và các ví dụ nằm trong văn bản đầu vào; mô hình được dùng
cho nhiệm vụ đó mà không cập nhật gradient hoặc fine-tune riêng cho nó
~\autocite{p03gpt3}. Vì vậy, thay một ví dụ trong lời nhắc có thể làm đổi câu
trả lời mà không hề thay đổi trọng số.

Điều này tạo ra hai câu hỏi giải thích khác nhau. Câu hỏi thứ nhất là token
nào trong lời nhắc làm thay đổi đầu ra đang xét. Câu hỏi thứ hai là mô hình
đã dùng tri thức nào trong trọng số để đáp lại token ấy. Các phương pháp ở
Phần~II thường tiếp cận câu hỏi thứ nhất vì ta có thể can thiệp trực tiếp vào
đầu vào. Câu hỏi thứ hai đòi hỏi bằng chứng về biểu diễn và cơ chế, nên không
nên được trả lời bằng một thay đổi đầu ra đơn lẻ.

\index{tiền huấn luyện}%
\index{lời nhắc}%
